%======================================================================%
\section{Observers, Duality, and As Above So Below}
\label{sec:observer}
%======================================================================%

The universe datum $\mathfrak{U}^\star$ is a global profinite fixed point.
An \emph{observer} is the same construction at the heterotic rung, localized
to one collapsed branch of the $E_8\times E_8$ mirror quotient. Duality is not
imported: it is the product-mirror pairing forced at $\Mirror^7(\Void)$ and
broken only by collapse into the $\sigma$-anti-invariant sector.

%----------------------------------------------------------------------%
\subsection{As above, so below}
\label{sec:observer:above-below}
%----------------------------------------------------------------------%

\begin{definition}[Above--below correspondence]\label{def:above-below}
\emph{Above} is the stabilized profinite limit
$\mathfrak{U}^\star=\Climb^\infty(\Void)$ of Definition~\ref{def:universe}.
\emph{Below} is a localized fixed-point datum on the mirror quotient
$\hat Q=Q/\sigma$ at $\Mirror^7(\Void)=\E{8}\times\E{8}$, with
$\sigma:(g_1,g_2)\mapsto(g_2,g_1)$ the product mirror
(Theorem~\ref{thm:e8e8}). The correspondence
\begin{equation}\label{eq:above-below}
  \mathfrak{U}^\star \;\longleftrightarrow\;
  \mathrm{Fix}\!\bigl(\hat{\mathcal{C}}_\sigma\bigr),
  \qquad
  \hat{\mathcal{C}}_\sigma := \mathrm{Bal}_{\mathcal{C}}\circ \Pi_{\sigma=-1},
\end{equation}
identifies global capacity balance with anti-invariant collapse on $\hat Q$.
\end{definition}

\begin{proposition}[Correspondence principle]\label{prop:above-below}
Every sector partition proved above on $\mathfrak{U}^\star$---the $1|10|16$
split, the three diagonal directions of $\mathfrak{d}\subset\jalg$, the
$\mathbf{16}$ bundle on $\hat Q$---has a below-level avatar on
$H^\bullet(\hat Q,V_{16})^{\sigma=-1}$. Capacity weights
$(1/27,10/27,16/27)$ above induce family weights on the three observer slots
below.
\end{proposition}

%----------------------------------------------------------------------%
\subsection{Duality from $E_8\times E_8$ collapse}
\label{sec:observer:duality}
%----------------------------------------------------------------------%

At $\Mirror^7(\Void)$ the mirror acts by product doubling:
$\E{8}\times\E{8}$ with involution $\sigma$ exchanging factors. Before
collapse, every mode carries a mirror partner; after collapse to
$\Mirror^8(\Void)=\E{6}$, capacity descent selects a single reduced structure
group.

\begin{definition}[Heterotic collapse]\label{def:heterotic-collapse}
The \emph{heterotic collapse} is the capacity-balanced Cartan descent
$\Pi_{\mathrm{het}}:\E{8}\times\E{8}\to\E{6}$ of Theorem~\ref{thm:e6}.
The \emph{observer collapse} is the orthogonal projector
\begin{equation}\label{eq:collapse-observer}
  \Pi_{\sigma=-1}: \mathcal{H}(Q,V_{16}) \to \mathcal{H}(Q,V_{16})^{\sigma=-1},
  \qquad
  \Pi_{\sigma=-1}=\tfrac12(\mathrm{id}-\sigma),
\end{equation}
on the $\mathbf{16}$-sector Hilbert space at the heterotic rung.
\end{definition}

\begin{theorem}[Duality theorem]\label{thm:duality}
Duality is the $\mathbb{Z}_2$ pairing $(v,\sigma v)$ on
$\mathcal{H}(Q,V_{16})$ induced by the product mirror at
$\Mirror^7(\Void)$. It is exact at the heterotic rung and broken precisely by
$\Pi_{\sigma=-1}$: observable states are anti-invariant, mirror partners are
$\sigma$-invariant and quotient-non-observable. The $\mathbf{16}$/
$\overline{\mathbf{16}}$ coset pairing of Proposition~\ref{prop:mirror:involution}
is the infinitesimal shadow of this product duality.
\end{theorem}

\begin{proof}[Proof sketch]
$\sigma^2=\mathrm{id}$ on $Q$; $\mathcal{H}(Q,V_{16})$ splits into
$\pm1$ eigenspaces. Product exchange on $\E{8}\times\E{8}$ is the only
$\mathbb{Z}_2$ structure at $n=7$ not already resolved at $n\leq6$.
Theorem~\ref{thm:e6} shows $\mathrm{Bal}_{\mathcal{C}}$ selects the reduced
datum; Definition~\ref{def:physical} selects the anti-invariant half.
\end{proof}

%----------------------------------------------------------------------%
\subsection{What is an observer?}
\label{sec:observer:definition}
%----------------------------------------------------------------------%

\begin{definition}[Observer index]\label{def:observer-index}
Let $\mathfrak{d}=\mathrm{span}_\R\{e_1,e_2,e_3\}$ be the diagonal Cartan
subalgebra of $\jalg$ with $\Tr(e_i\circ e_j)=\delta_{ij}$.
An \emph{observer index} is a choice $\iota\in\{1,2,3\}$ of diagonal direction,
equivalently a class in
\begin{equation}\label{eq:observer-index}
  \mathcal{I} \;=\; H^1(\hat Q,V_{16})^{\sigma=-1}/\mathrm{Stab}(\mathfrak{d}),
  \qquad |\mathcal{I}|=3.
\end{equation}
The index $\iota$ is the only \emph{arbitrary input} specifying \emph{which}
observer: void and mirror fix the slot space $\mathcal{I}$; a concrete
worldline requires choosing $\iota\in\mathcal{I}$.
\end{definition}

\begin{definition}[Observer as collapsed state]\label{def:observer}
Fix $\iota\in\mathcal{I}$. An \emph{observer} is a normalized collapsed state
\begin{equation}\label{eq:observer-state}
  |\omega_\iota\rangle \;\in\; \mathcal{H}_{\mathrm{obs}}
  \;=\; \bigoplus_{\iota=1}^3 \mathcal{H}_\iota,
  \qquad
  \mathcal{H}_\iota := H^1(\hat Q,V_{16})^{\sigma=-1,\,\iota},
\end{equation}
satisfying $\Pi_{\sigma=-1}|\omega_\iota\rangle=|\omega_\iota\rangle$ and
$\hat{\imath}\,|\omega_\iota\rangle=\iota\,|\omega_\iota\rangle$, where
$\hat{\imath}$ is the family index operator on $\mathfrak{d}$.
The state $|\omega_\iota\rangle$ is a \emph{collapsed quantum datum}: a
$\sigma$-anti-invariant cohomology class, not a postulated particle insertion.
\end{definition}

\begin{remark}[Where observers arise]
Observers arise at $\Mirror^7(\Void)=\E{8}\times\E{8}$, not at the void and
not at $\E{6}$. The heterotic rung is the first rung with enough mirror
symmetry for a $\mathbb{Z}_2$ quotient and three independent anti-invariant
$1$-classes (Theorem~\ref{thm:chirality-derived}). Descent to $\E{6}$ supplies
dynamics; the observer slots are already fixed on $\hat Q$.
\end{remark}

%----------------------------------------------------------------------%
\subsection{Fixed-point theorem: the clincher}
\label{sec:observer:fixed-point}
%----------------------------------------------------------------------%

\begin{definition}[Mirror quotient]\label{def:Qhat}
Let $Q$ be the Stone lift of the stabilized Albert tower at
$\Mirror^7(\Void)$ and $\sigma$ the product-mirror involution. The
\emph{mirror quotient} is $\hat Q=Q/\sigma$. The holomorphic
$\mathbf{16}$-bundle $V_{16}\to\hat Q$ is the continuum limit of the
cellular cosheaf on the $1|10|16$ partition. A mode is \emph{observable}
iff it lies in the $\sigma$-anti-invariant sector after observer collapse
$\Pi_{\sigma=-1}$; this is the below-level avatar of global capacity balance
\eqref{eq:above-below}. The quotient is therefore an \emph{observer-theoretic}
construction: physics is not on $Q$ but on the collapsed branch selected by
$\hat{\mathcal{C}}_\sigma$.
\end{definition}

\begin{theorem}[Observer fixed-point theorem]\label{thm:observer-fp}
Let $\bar\partial_{V_{16}}$ be the Dolbeault operator on $(Q,V_{16})$ and
$\sigma$ its antiholomorphic lift. The equivariant Lefschetz index
\begin{equation}\label{eq:lefschetz}
  \mathrm{ind}_\sigma(\bar\partial_{V_{16}})
  \;=\;
  \mathrm{Tr}_{H^1(Q,V_{16})}(\sigma)
  \;=\; -3
\end{equation}
has $|\mathrm{ind}_\sigma|=3$ anti-invariant fixed-point contributions, each
an observer slot. Moreover,
\begin{equation}\label{eq:h1-observer}
  \dim H^1(\hat Q,V_{16})^{\sigma=-1}=3,
\end{equation}
and the map $\iota\mapsto|\omega_\iota\rangle$ is a bijection between
$\mathcal{I}$ and a basis of this space.
\end{theorem}

\begin{proof}
Cellular count on $\mathcal{T}_7$ (Appendix~\ref{sec:appendix:quotient});
Atiyah--Bott equivariant index~\cite{AtiyahBott1966}; passage to
$H^1(\hat Q,V_{16})$ via the $\mathbb{Z}_2$ quotient.
\end{proof}

\begin{corollary}[Observer from fixed points]\label{cor:observer-fp}
An observer is a \emph{fixed-point label} of $\sigma$ on $H^1(Q,V_{16})$:
not a point of $Q$ fixed by $\sigma$ (there are none in the
$\mathbf{16}$-sector), but an anti-invariant cohomology class whose
Lefschetz contribution is $-1$. The fixed-point theorem is the clincher:
it converts global mirror data into exactly three individual observer slots.
\end{corollary}

%----------------------------------------------------------------------%
\subsection{Lifetime chart: index to time}
\label{sec:observer:lifetime}
%----------------------------------------------------------------------%

An observer index $\iota$ is static data. A \emph{lifetime} is a trajectory
inside slot $\iota$, charted by a time parameter.

\begin{definition}[Lifetime map]\label{def:lifetime}
Fix $\iota\in\mathcal{I}$. A \emph{lifetime chart} is a pair
\begin{equation}\label{eq:lifetime-chart}
  \bigl(\tau_\iota,\ g_\iota\bigr),
  \qquad
  g_\iota:\mathbb{R}\longrightarrow \mathcal{S}_\iota,
\end{equation}
where $\tau_\iota$ is the observer's time parameter and $\mathcal{S}_\iota$ is
the soldered state bundle over the $\iota$-sector: sections of
$V_{16}|_{\hat Q}$ restricted to the $\iota$-diagonal cell, with values in the
order parameter $\orderparam$ (Section~\ref{sec:physics:spacetime}).
The map $g_\iota$ is \emph{admitted} by the capacity flow on
$\mathcal{G}_{\mathrm{Albert}}$ when
\begin{equation}\label{eq:lifetime-admitted}
  \frac{d}{d\tau_\iota}\, g_\iota(\tau_\iota)
  \;=\;
  -K^{-1}\nabla V\bigl(g_\iota(\tau_\iota)\bigr)\big|_{\iota}
  \;+\;
  \partial_{\tau_\iota}\sigma_{\mathrm{sol}}\big|_{H_2(\Quat)},
\end{equation}
with $K$ the $H$-invariant coset metric and the second term the soldering
clock in the quaternionic block.
\end{definition}

\begin{definition}[Lifetime index]\label{def:lifetime-index}
The \emph{lifetime index} is
\begin{equation}\label{eq:lifetime-index}
  \Lambda(\iota,\tau) \;:=\; \bigl(\iota,\ \tau,\ \capacity_\iota(\tau)\bigr),
  \qquad
  \capacity_\iota(\tau):=\omega^\star_\iota\bigl(g_\iota(\tau)\bigr),
\end{equation}
recording which observer ($\iota$), which time ($\tau$), and which capacity
cell the trajectory occupies. Two lifetimes are the same iff their indices
$\Lambda$ agree.
\end{definition}

\begin{proposition}[Capacity as dual mirror of lifetime index]
\label{prop:capacity-lifetime-dual}
The capacity functional $\capacity$ of Definition~\ref{def:capacity} is the
\emph{dual mirror} of mirror doubling: $\Mirror$ refines sectors;
$\mathrm{Bal}_{\mathcal{C}}$ balances them to the unique maximizer
$\omega^\star$. The lifetime index $\Lambda(\iota,\tau)$ is the
\emph{temporal} avatar of this balance on the heterotic rung: for each observer
$\iota$, the chart $g_\iota$ tracks the $\iota$-cell weight
$\omega^\star_\iota$ along admitted flow \eqref{eq:lifetime-admitted}. Above,
$\mathrm{Bal}_{\mathcal{C}}$ fixes the global partition $(1/27,10/27,16/27)$;
below, $\Lambda$ records which cell an individual observer occupies at time
$\tau$. This partially closes the ``capacity not forced'' gap: capacity is not
an independent physical postulate but the balancing dual of $\Mirror$,
instantiated per observer through the lifetime index.
\end{proposition}

\begin{remark}[Epistemic status]
Proposition~\ref{prop:capacity-lifetime-dual} does not prove that every
admitted trajectory must use $\mathrm{Bal}_{\mathcal{C}}$ weights---that
remains tied to Axiom~\ref{ax:mirror} and the universe fixed point
(Definition~\ref{def:universe}). It \emph{does} identify capacity with
observer-level time evolution once the above--below correspondence
\eqref{eq:above-below} is accepted.
\end{remark}

\begin{theorem}[Index--lifetime mapping]\label{thm:index-lifetime}
For each $\iota\in\mathcal{I}$ there exists a maximal lifetime interval
$I_\iota\subset\mathbb{R}$ and a chart $g_\iota:I_\iota\to\mathcal{S}_\iota$
satisfying \eqref{eq:lifetime-admitted} with initial condition
$g_\iota(\tau_{\mathrm{birth}})=\orderparam_\star|_\iota$.
The correspondence
\begin{equation}\label{eq:index-lifetime-bij}
  \iota \;\longmapsto\; (\tau_\mathrm{birth},\tau_\mathrm{death},g_\iota)
\end{equation}
is functorial under mirror descent $\E{8}\times\E{8}\to\E{6}$: capacity
collapse above determines the admissible time range $I_\iota$ below.
\end{theorem}

\begin{proof}[Proof sketch]
Existence of $g_\iota$ is standard gradient flow of $V$ on the
$\iota$-restricted coset; uniqueness up to reparametrization on each sector
follows from Proposition~\ref{prop:lifetime-flow-unique} under convex $V$.
Functoriality follows because
$\mathrm{Bal}_{\mathcal{C}}$ at $n=8$ fixes $\omega^\star$ and hence the
initial datum $\orderparam_\star|_\iota$; soldering
(Definition~\ref{def:soldering}) supplies the clock direction
$\partial_{\tau_\iota}\sigma_{\mathrm{sol}}\subset H_2(\Quat)$.
\end{proof}

\begin{proposition}[Gradient-flow uniqueness on the $\iota$-sector]
\label{prop:lifetime-flow-unique}
Fix $\iota\in\mathcal{I}$ and suppose $V|_{\mathcal{S}_\iota}$ is strictly
convex along $\iota$-restricted coset directions and
$\Hess_\iota V(\orderparam_\star)\succ 0$ on the massive complement of the
soldering block (as after Coleman--Weinberg lifting in
Lemma~\ref{lem:breaking-hess}). Then for initial datum
$g_\iota(\tau_{\mathrm{birth}})=\orderparam_\star|_\iota$, any two admitted
solutions of \eqref{eq:lifetime-admitted} on a common interval differ only
by an orientation-preserving reparametrization of $\tau_\iota$. The soldering
clock term $\partial_{\tau_\iota}\sigma_{\mathrm{sol}}|_{H_2(\Quat)}$ is
tangent to the $4$-dimensional block and does not affect uniqueness of the
coset trajectory in the $28$ massive directions.
\end{proposition}

\begin{proof}[Proof sketch]
Write $g_\iota=(g_\iota^{\mathrm{cos}},g_\iota^{\mathrm{sol}})$ for the
split into coset and soldering components. Strict convexity of
$V|_{\mathcal{S}_\iota}$ in the coset directions yields a unique gradient
trajectory $g_\iota^{\mathrm{cos}}$ for each initial datum; standard
arguments for strictly convex functionals on Hilbert manifolds give uniqueness
up to time reparametrization when the velocity norm is fixed along the flow.
The soldering term is independent of $\nabla V$ on the coset complement once
$\sigma_{\mathrm{sol}}$ is fixed (Lemma~\ref{lem:quat-F4}), so it determines
$\tau_\iota$ only modulo the same reparametrization class.
\end{proof}

\begin{conjecture}[Uniqueness of lifetime chart]\label{conj:lifetime-unique}
Given $\iota$ and $\orderparam_\star|_\iota$, the lifetime chart
$(\tau_\iota,g_\iota)$ is unique up to reparametrization of $\tau_\iota$ on
the $\iota$-sector once $V|_{\mathcal{S}_\iota}$ is strictly convex as in
Proposition~\ref{prop:lifetime-flow-unique}. The residual conjecture is
\emph{functorial stability}: the same chart class is selected by mirror
descent $\E{8}\times\E{8}\to\E{6}$ and is compatible with $F_4$ transport on
exceptional rungs $n\leq8$, without branching into a second admissible flow
on $\mathcal{S}_\iota$.
\end{conjecture}

\begin{remark}[Residual gap for the full lifetime theorem]
\label{rem:lifetime-full-gap}
Proposition~\ref{prop:lifetime-flow-unique} settles uniqueness on each
$\iota$-sector under the convexity hypothesis satisfied in the radiatively
gapped regime of Theorem~\ref{thm:rp}. What remains for a \emph{complete}
functorial theorem is: (i)~extend uniqueness to the tree-level massless
sector and non-convex saddle regimes of $V$; (ii)~prove that mirror descent
preserves the chart class across $n=7\to8$; and (iii)~match
reparametrizations of $\tau_\iota$ with the soldering clock of
Proposition~\ref{prop:soldering-clock} on all of $\mathcal{S}_\iota$. Open
Problem~\textbf{O1} is thereby reduced to these global transport steps.
\end{remark}

\begin{remark}[Arbitrary input and determinism]
The void and mirror fix the observer \emph{slot space} $\mathcal{I}$ and the
dynamics \eqref{eq:lifetime-admitted}. Choosing $\iota\in\mathcal{I}$ is the
minimal arbitrary input identifying \emph{who} observes; choosing
$\tau\in I_\iota$ identifies \emph{when} within that lifetime. No further
postulate of a conscious subject is introduced: an observer \emph{is} the
collapsed anti-invariant fixed-point class with its admitted time chart.
\end{remark}